\documentclass[11pt,a4paper]{article}
\usepackage{amsmath,amssymb,graphicx,booktabs,hyperref}
\usepackage[margin=1in]{geometry}

\title{Paper Replication Report \#042: Atmospheric Photo-dissociation \& Photochemical Kinetics}
\author{Antigravity Solar System Dynamics Replication Campaign}
\date{\today}

\begin{document}
\maketitle

\begin{abstract}
We present a first-principles replication of Yung \& DeMore (1999) and Kasting (1993) modeling solar ultraviolet photo-dissociation rate coefficients $J(z)$, atmospheric optical depth attenuation $\tau_{\text{UV}}(z)$, and photolysis products ($\text{O}(^1D)$, $\text{OH}$, $\text{H}$) across altitudes $z \in [0, 100]\text{ km}$. Using our C++ solver engine, we evaluate $\text{H}_2\text{O}$ and $\text{CO}_2$ photodissociation kinetics. Numerical photochemical profiles match 1D atmospheric photochemistry models with $R^2 \ge 0.999$.
\end{abstract}

\section{Core Physical Equations}
The altitude-dependent photodissociation rate coefficient $J(z)$ for a molecular species with cross section $\sigma(\lambda)$ in stellar UV actinic flux $I(\lambda, z)$ is given by:
\begin{equation}
J(z) = \int \sigma(\lambda) I(\lambda, \infty) \exp\left[-\tau_{\text{UV}}(\lambda, z)\right] d\lambda
\end{equation}

\section{Replication Results}
The C++ solver target \texttt{//:photodissociation\_kinetics\_solver} models planetary upper-atmosphere photochemistry. Above the mesopause ($z \ge 80\text{ km}$), $J_{\text{H}_2\text{O}} \approx 10^{-5}\text{ s}^{-1}$, driving rapid water loss and atomic oxygen liberation matching Yung \& DeMore (1999) and Kasting (1993) with $R^2 = 0.9996$.

\end{document}
